\section*{Capítulo \ref{ch:g6:circuits-and-safety} --- Circuitos elétricos e segurança}

\begin{solution}{exo:g6:circuits-and-safety:1}
Com fidelidade: as ligações --- qual dispositivo se junta a qual,
onde as estradas se dividem. Ignorados: o comprimento, a cor e o
serpentear dos fios de verdade sobre a mesa.
\end{solution}

\begin{solution}{exo:g6:circuits-and-safety:2}
\omterm{def:g3:first-electric-circuit:battery}{Pilha}: duas barras paralelas de tamanhos diferentes; lâmpada: círculo
com uma cruz; \omterm{ex:g3:first-electric-circuit:switch}{interruptor} aberto: uma ponte levadiça levantada;
\omterm{ex:g3:first-electric-circuit:switch}{interruptor} fechado: a ponte abaixada; \omterm{ex:g6:circuits-and-safety:motor}{motor}: um M num círculo;
\omterm{ex:g6:circuits-and-safety:motor}{campainha}: um círculo com um traço. A barra \emph{comprida} marca o
$+$ da \omterm{def:g3:first-electric-circuit:battery}{pilha}.
\end{solution}

\begin{solution}{exo:g6:circuits-and-safety:3}
Um retângulo: \omterm{def:g3:first-electric-circuit:battery}{pilha}, botão e \omterm{ex:g6:circuits-and-safety:motor}{campainha} numa única volta. Tirar o dedo
abre o \omterm{ex:g3:first-electric-circuit:switch}{interruptor} --- uma falha na única estrada --- e a lei da
volta silencia a \omterm{ex:g6:circuits-and-safety:motor}{campainha} na hora.
\end{solution}

\begin{solution}{exo:g6:circuits-and-safety:4}
Não necessariamente. A pergunta que decide: quem está ligado a quem?
Os mesmos dispositivos com as mesmas ligações --- mesmo circuito, por
mais diferentes que os emaranhados pareçam.
\end{solution}

\begin{solution}{exo:g6:circuits-and-safety:5}
O ventilador continua girando: a morte da lâmpada apaga só o \omterm{def:g5:series-parallel-circuits:parallel}{ramo}
dela --- é a independência do paralelo. O \omterm{ex:g3:first-electric-circuit:switch}{interruptor} aberto fica na
estrada comum, antes da divisão, então ele para os dois: a regra dos
postos de comando.
\end{solution}

\begin{solution}{exo:g6:circuits-and-safety:6}
Uma estrada condutora nua ligando diretamente os dois lados da fonte,
pulando todos os dispositivos: a corrente dispara sem obstáculo e os
fios esquentam. Vítimas: a \omterm{def:g3:first-electric-circuit:battery}{pilha} (descarregada e arruinada), o
isolamento do fio --- e, na rede elétrica, às vezes a casa inteira.
\end{solution}

\begin{solution}{exo:g6:circuits-and-safety:7}
Ele cortou o \omterm{def:g5:series-parallel-circuits:parallel}{ramo} dele num piscar porque a corrente começou a
disparar --- um defeito, e não um capricho. Encontre e conserte o
defeito (tire da tomada o suspeito, inspecione o \omterm{def:g5:series-parallel-circuits:parallel}{ramo}) antes de
rearmar o guarda.
\end{solution}

\begin{solution}{exo:g6:circuits-and-safety:8}
Banheira: a pele úmida conduz, e o empurrão da rede através de um
corpo molhado pode ser mortal. Fios danificados: isolamento rachado é
uma cerca com buraco --- os dedos podem encontrar o cobre. \omterm{def:g3:first-electric-circuit:battery}{Pilha} sim,
tomada nunca: o empurrão da \omterm{def:g3:first-electric-circuit:battery}{pilha} é pequenininho; o da tomada é
centenas de vezes mais possante.
\end{solution}

\begin{solution}{exo:g6:circuits-and-safety:9}
Esquema: \omterm{def:g3:first-electric-circuit:battery}{pilha}, \omterm{ex:g3:first-electric-circuit:switch}{interruptor} fechado na estrada comum e depois uma
divisão --- \omterm{def:g5:series-parallel-circuits:parallel}{ramo} da \omterm{ex:g6:circuits-and-safety:motor}{campainha} e \omterm{def:g5:series-parallel-circuits:parallel}{ramo} da lâmpada --- juntando-se de
volta até o $-$. Previsão: a lâmpada acende \emph{e} a \omterm{ex:g6:circuits-and-safety:motor}{campainha} soa,
de forma independente. Com um \omterm{ex:g3:first-electric-circuit:switch}{interruptor} extra aberto no \omterm{def:g5:series-parallel-circuits:parallel}{ramo} da
\omterm{ex:g6:circuits-and-safety:motor}{campainha}: silêncio ali, mas a lâmpada continua acesa --- os
\omterm{ex:g3:first-electric-circuit:switch}{interruptores} de \omterm{def:g5:series-parallel-circuits:parallel}{ramo} mandam só no \omterm{def:g5:series-parallel-circuits:parallel}{ramo} deles.
\end{solution}

\begin{solution}{exo:g6:circuits-and-safety:10}
O fio extra é um \omterm{def:g6:circuits-and-safety:short}{curto-circuito}: uma estrada nua e direta do $+$ ao
$-$. A corrente prefere a estrada fácil e vazia ao \omterm{ex:g6:circuits-and-safety:motor}{motor} trabalhador,
então o ventilador passa fome enquanto o fio-atalho carrega a
disparada e esquenta --- descarregando a \omterm{def:g3:first-electric-circuit:battery}{pilha} e convidando
queimaduras. Tire o fio da ``firmeza''.
\end{solution}

\begin{solution}{exo:g6:circuits-and-safety:11}
O \omterm{ex:g3:first-electric-circuit:switch}{interruptor} de parede abre a volta da própria luminária --- a
estrada da lâmpada está cortada. O disjuntor abre o \omterm{def:g5:series-parallel-circuits:parallel}{ramo} inteiro mais
acima --- assim, mesmo que a fiação da luminária esteja com defeito
ou o \omterm{ex:g3:first-electric-circuit:switch}{interruptor} de parede esteja mal identificado, o \omterm{def:g5:series-parallel-circuits:parallel}{ramo} está
morto. Duas falhas \omterm{def:g4:series-circuits:series}{em série}: qualquer uma sozinha já para a corrente,
e as duas juntas perdoam um engano sobre qualquer uma delas.
\end{solution}

\begin{solution}{exo:g6:circuits-and-safety:12}
\omterm{def:g3:first-electric-circuit:battery}{Pilha}, depois o \omterm{ex:g3:first-electric-circuit:switch}{interruptor} geral na estrada comum e depois quatro
\omterm{def:g5:series-parallel-circuits:parallel}{ramos} paralelos: \omterm{def:g1:light-and-shadows:source}{luz} de palco 1; \omterm{def:g1:light-and-shadows:source}{luz} de palco 2; \omterm{ex:g6:circuits-and-safety:motor}{motor} da cortina com
o \omterm{ex:g3:first-electric-circuit:switch}{interruptor} \omterm{def:g4:series-circuits:series}{em série} dele; \omterm{ex:g6:circuits-and-safety:motor}{campainha} com o \omterm{ex:g3:first-electric-circuit:switch}{interruptor} \omterm{def:g4:series-circuits:series}{em série}
dela. Conferências: o geral manda em tudo (estrada comum); as \omterm{def:g1:light-and-shadows:source}{luzes}
falham de forma independente (\omterm{def:g5:series-parallel-circuits:parallel}{ramos} separados); o \omterm{ex:g6:circuits-and-safety:motor}{motor} e a \omterm{ex:g6:circuits-and-safety:motor}{campainha}
obedecem cada um ao seu \omterm{ex:g3:first-electric-circuit:switch}{interruptor} (\omterm{def:g4:series-circuits:series}{em série} dentro do \omterm{def:g5:series-parallel-circuits:parallel}{ramo}) e
ignoram um ao outro (\omterm{def:g5:series-parallel-circuits:parallel}{em paralelo} entre \omterm{def:g5:series-parallel-circuits:parallel}{ramos}).
\end{solution}

\begin{solution}{pb:g6:circuits-and-safety:1}
\textbf{1.} \omterm{def:g4:series-circuits:series}{Em série}, um dispositivo queimado ou desligado apagaria
tudo, e os três dispositivos repartiriam o empurrão, todos fracos e
mortiços. O paralelo mantém cada dispositivo independente e na \omterm{def:g2:pushes-and-pulls:force}{força}
total.
\textbf{2.} Cada \omterm{ex:g3:first-electric-circuit:switch}{interruptor} manda exatamente no dispositivo do
próprio \omterm{def:g5:series-parallel-circuits:parallel}{ramo}; nenhum deles manda nos outros \omterm{def:g5:series-parallel-circuits:parallel}{ramos}.
\textbf{3.} Na estrada comum, entre a \omterm{def:g3:first-electric-circuit:battery}{pilha} e a primeira divisão ---
antes de todos os \omterm{def:g5:series-parallel-circuits:parallel}{ramos}.
\textbf{4.} \omterm{def:g3:first-electric-circuit:battery}{Pilha} $\to$ \omterm{ex:g3:first-electric-circuit:switch}{interruptor} geral $\to$ \omterm{def:g5:series-parallel-circuits:parallel}{nó} que se divide em
três \omterm{def:g5:series-parallel-circuits:parallel}{ramos} --- (luminária de teto $+$ \omterm{ex:g3:first-electric-circuit:switch}{interruptor} dela), (luminária
de bancada $+$ \omterm{ex:g3:first-electric-circuit:switch}{interruptor} dela), (\omterm{ex:g6:circuits-and-safety:motor}{motor} do ventilador $+$
\omterm{ex:g3:first-electric-circuit:switch}{interruptor} dele) --- juntando-se de novo e voltando à \omterm{def:g3:first-electric-circuit:battery}{pilha}.
\textbf{5.} Eles estão \omterm{def:g4:series-circuits:series}{em série} num mesmo \omterm{def:g5:series-parallel-circuits:parallel}{ramo}: repartem o empurrão
(os dois fracos --- a luminária de teto fraquinha era, na verdade, o
sintoma da dupla da bancada), e o \omterm{ex:g3:first-electric-circuit:switch}{interruptor} do ventilador, que fica
na estrada compartilhada, manda nos dois.
\textbf{6.} Separe-os: passe um fio novo para que a luminária de
bancada ganhe o próprio \omterm{def:g5:series-parallel-circuits:parallel}{ramo} ligado à alimentação --- assim cada uma
fica na \omterm{def:g2:pushes-and-pulls:force}{força} total e independente.
\textbf{7.} Um \omterm{def:g6:circuits-and-safety:short}{curto-circuito} entre os \omterm{def:g3:first-electric-circuit:battery}{terminais}. Confira se a \omterm{def:g3:first-electric-circuit:battery}{pilha}
ainda entrega corrente --- os curtos descarregam e podem arruinar as
células (uma \omterm{def:g3:first-electric-circuit:battery}{pilha} morna merece aposentadoria).
\textbf{8.} Curtos-circuitos (fios nus se encontrando) e choques ou
queimaduras (dedos encontrando cobre nu).
\textbf{9.} A pele úmida conduz muito melhor que a seca; até a
montagem com \omterm{def:g3:first-electric-circuit:battery}{pilha} merece a disciplina das mãos secas --- e, na casa,
o mesmo toque encontra o empurrão da rede, centenas de vezes mais
possante: possivelmente mortal. Os hábitos precisam ser treinados no
circuito seguro.
\textbf{10.} O disjuntor detectou que corrente demais estava
disparando pelo \omterm{def:g5:series-parallel-circuits:parallel}{ramo} dele --- um defeito, provavelmente chuva numa
luminária. Ordem sábia: deixar o disjuntor desligado, desligar e
inspecionar o \omterm{def:g5:series-parallel-circuits:parallel}{ramo}, consertar ou secar o defeito e só então rearmar.
\textbf{11.} Uma volta própria, separada da rede do galpão (ou como
mais um \omterm{def:g5:series-parallel-circuits:parallel}{ramo} \omterm{def:g5:series-parallel-circuits:parallel}{em paralelo} na \omterm{def:g3:first-electric-circuit:battery}{pilha} do galpão): \omterm{def:g3:first-electric-circuit:battery}{pilha}, botão na casa,
fio duplo comprido, \omterm{ex:g6:circuits-and-safety:motor}{campainha} no galpão e de volta --- uma \omterm{ex:g6:circuits-and-safety:motor}{campainha}
de porta é uma volta \omterm{def:g4:series-circuits:series}{em série} própria, por mais longe que os fios
dela se estiquem.
\textbf{12.} Por exemplo: ``Tomadas: as nossas experiências terminam
na \omterm{def:g3:first-electric-circuit:battery}{pilha} --- a tomada não é nossa. Água: mãos molhadas não tocam
fiação nenhuma, dentro nem fora de casa. Fios nus: toda emenda
encapada; um fio nu é um defeito, e não um atalho.''
\end{solution}
